% ============================================================
%  Network mode
%  Multi-station PPP-RTK network estimator that GENERATES the
%  SSR-style corrections the single-receiver modes consume:
%  joint observation model, S-basis / minimum-constraint datum,
%  per-station Schur reduction, explicit satellite phase-bias
%  state.
%
%  Sources: filter/NetworkEstimator + NetworkNormalEquations
%  (the arrowhead Schur kernel), filter/StateManager (explicit
%  sat-phase-bias block), measurement/MeasurementBuilder
%  (DifferencingMode::None), app/ApplicationNetwork.cpp, and the
%  S-system theory in foundations (\Cref{sec:diff}).
% ============================================================
\documentclass[main.tex]{subfiles}
\begin{document}
\ifSubfilesClassLoaded{\begin{center}\small\itshape Part of the GnssProcessor Mathematical Reference \textbullet\ \today\end{center}}{}

\section{Network mode}\label{sec:net}

\newcommand{\Nss}{\bm{N}_{ss}}          % shared (satellite) normal block
\newcommand{\Nrr}{\bm{N}_{rr}}          % per-station local normal block
\newcommand{\Nsr}{\bm{N}_{sr}}          % shared<->local cross block
\newcommand{\redNss}{\tilde{\bm{N}}_{ss}} % Schur-reduced shared normals
\newcommand{\redus}{\tilde{\bm{u}}_s}   % Schur-reduced shared RHS
\newcommand{\xshared}{\bm{x}_s}         % shared satellite-parameter estimate
\newcommand{\xlocal}[1]{\bm{x}_{#1}}    % station r's private nuisance block
\newcommand{\satbiasPB}{\bar b}         % satellite phase bias (estimable); matches \bar b in gnssreport.sty
\newcommand{\satbiasCB}{B}              % satellite code bias (estimable)

This section documents the multi-station network estimator
(\texttt{ProcessingMode::PPP\_RTK\_Network}) that \emph{produces} the
satellite-side corrections the single-receiver user modes consume.
Where every other mode in this reference \emph{differences away} the
satellite clock and hardware biases
(\dcref{sec:diff}{the differencing section}) or imports them
as known SSR quantities (\dcref{sec:ssr}{the SSR-corrections section}),
the network producer turns that relationship around: it estimates the
satellite clock and the satellite \emph{phase} and \emph{code} biases
\emph{explicitly}, jointly across a network of stations whose
coordinates are known, and emits them in IGS-SSR convention for a
downstream consumer. The closed loop is deliberate --- the same
single-receiver PPP-RTK engine documented elsewhere becomes the
product test, run on a station \emph{held out} of the network solve.

\paragraph{Scope and status.}
The numerical core --- the arrowhead normal-equation kernel
(\texttt{filter/NetworkNormalEquations}) and the semantic estimator
that drives it (\texttt{filter/NetworkEstimator}) --- is implemented
and validated on a synthetic noise-free network (truth recovery up to
the declared S-transform). The full pipeline is live: the
\texttt{PPP\_RTK\_Network} enum, its forced-uncombined
\texttt{DifferencingMode::None}
(\texttt{src/core/measurement/MeasurementBuilder.cpp}), the explicit
satellite phase-bias state in \texttt{filter/StateManager}, and the
dispatch in \texttt{Application::run} all exist, and the epoch loop
\texttt{Application::runNetwork} is fully wired: per-station design-row
assembly from each station's \texttt{MeasurementModel}, the S-basis
datum with its null-space rank gate (\Cref{sec:net-sbasis}), the
between-epoch time update, and the product emitter
(\texttt{NetworkProductEmitter}) all run. The mathematics below
describes the implemented estimator core and the locked design. This
chapter treats the single-epoch solve; the between-epoch carry
machinery --- the C1 full-covariance shared carry
(\texttt{addSharedPriorDense}), the Card-4 invariant DD-carry
(\texttt{setDdPriors}), and the MW-WL / EWL / lane-coordinate widelane
layers --- is a large, default-off subsystem in its own right and is
only referenced here, not derived.

\subsection{Joint multi-station observation model}\label{sec:net-model}

The producer runs on the same \emph{undifferenced, uncombined}
observations as single-receiver PPP
(\dcref{sec:obs}{the observation-equations section}),
but stacks the equations of every station $\rx \in \rxset$ into one
system. The defining structural fact is the partition of the unknowns
into a block \emph{shared} by all stations and a block \emph{private}
to each:
\begin{itemize}[leftmargin=*, itemsep=1pt]
\item \textbf{Shared (satellite-side).} The estimable satellite clock
      $\sLight\,\mathrm{d}t^\sat$, the per-signal satellite code bias
      $\satbiasCB^\sat_\sigl$, and the per-signal satellite \emph{phase}
      bias $\satbiasPB^\sat_\sigl$. These columns are \emph{identical}
      across every station --- a satellite's clock and hardware delays
      do not depend on who is observing it. (Orbits are held fixed to a
      precise product here; orbit determination is out of scope for the
      producer.) When \texttt{network.estimate\_iono\_sh} is set, the
      shared block carries a fourth family: the coefficients
      $a_{nm},b_{nm}$ of a global single-layer spherical-harmonic VTEC
      expansion (\texttt{SharedParamKind::VtecSh}, registered through
      \texttt{NetworkEstimator::registerVtecCoeff}), network-wide rather
      than per-satellite. This co-estimated ionosphere model replaces
      the external GIM and forces satellite code-bias estimation on ---
      the two are jointly load-bearing for the per-satellite slant-iono
      columns below --- and its columns Schur-reduce onto the shared
      border exactly like the satellite-indexed families.
\item \textbf{Private (station-side).} For each station $\rx$: its
      receiver clock $\sLight\,\mathrm{d}t_\rx$, the inter-system time
      offsets $\sLight\,\delta t^{(g)}_\rx$, the zenith wet troposphere
      $\zwd_\rx$ (and gradients), the per-satellite slant ionosphere
      $I^\sat_{\rx,1}$, the receiver biases, and one float ambiguity
      $N^\sat_{\rx,\sigl}$ per tracked carrier. Station coordinates are
      \emph{known} (the network constrains them), so position is not a
      free parameter.
\end{itemize}

Write the shared satellite-parameter vector as $\xshared$ and station
$\rx$'s private nuisance vector as $\xlocal{\rx}$. The undifferenced
observation rows of station $\rx$ linearise
(\dcref{sec:obs}{the observation-equations section}) to
\begin{equation}
\bm{y}_\rx \;=\; \bm{A}_{s,\rx}\,\xshared \;+\; \bm{A}_{r,\rx}\,\xlocal{\rx}
            \;+\; \bm{\varepsilon}_\rx,
\qquad \rx \in \rxset,
\label{eq:net-station-rows}
\end{equation}
where $\bm{y}_\rx$ is the prefit residual (observed minus computed),
$\bm{A}_{s,\rx}$ is the Jacobian with respect to the \emph{shared}
satellite columns, and $\bm{A}_{r,\rx}$ is the Jacobian with respect to
station $\rx$'s \emph{own} nuisances. Crucially there is no
$\bm{A}_{r,\rx'}$ term for $\rx' \neq \rx$: a station's observations
touch only the shared block and its own private block. That sparsity is
the structural engine of \Cref{sec:net-schur}.

\paragraph{No receiver-side differencing.}
Unlike the PPP-RTK user mode, the producer applies \emph{no}
between-satellite or between-receiver differencing. The differencing
strategy is forced to \texttt{DifferencingMode::None}
(\texttt{MeasurementBuilder.cpp}), overriding any
\texttt{between\_satellite\_sd} toggle that may be present in a network
config. Differencing would cancel exactly the satellite-side terms the
producer is trying to \emph{estimate}; the network instead carries them
as explicit columns and resolves the resulting rank deficiency with a
minimum-constraint datum (\Cref{sec:net-sbasis}).

\subsection{The explicit satellite phase-bias state}\label{sec:net-phasebias}

The single defining contrast with the user side is the treatment of the
satellite phase bias $b^\sat_\sigl$ (here written $\satbiasPB^\sat_\sigl$
once promoted to a state). On the user side this term is never an
explicit unknown: between-satellite single differencing
(\dcref{sec:diff}{the differencing section}) cancels the
\emph{receiver} phase bias, the satellite phase bias arrives as a
\emph{known} SSR product, and what remains is an integer
single-difference ambiguity. The user handles biases \emph{implicitly}.

The producer is what \emph{generates} that known product, so it must
carry the satellite phase bias as a genuine estimated state. The state
vector therefore gains a per-$(\text{satellite},\text{signal})$ phase
bias column,
\begin{equation}
\satbiasPB^\sat_\sigl \in \xshared,
\qquad
\text{keyed by } \texttt{makeSatPhaseBiasKey}(\sat,\sigl)
   = \texttt{"$\sat$|$\sigl$|PB"},
\label{eq:net-pb-key}
\end{equation}
held in \texttt{StateManager::sat\_phase\_bias\_indices\_}, a map kept
\emph{deliberately disjoint} from the satellite \emph{code} bias map
(\texttt{sat\_bias\_indices\_}); the trailing \texttt{|PB}
discriminator keeps the two namespaces from colliding
(\texttt{include/filter/StateKeySchema.hpp}). Its process noise is a
tiny random walk (\texttt{process\_noise.sat\_phase\_bias}, default two
orders of magnitude below the code-bias density), reflecting that
satellite phase hardware delays are nearly constant but not strictly
so: a $q=0$ random walk would over-converge the diagonal of $\covM$ and
leave the column deaf to genuine slow drift
(\texttt{src/core/filter/StateManager.cpp}).

\begin{remark}
The estimated $\satbiasPB^\sat_\sigl$ is not the ``raw'' satellite phase
hardware delay --- that quantity is not separately estimable. It is the
delay \emph{lumped with the network pivot-receiver ambiguity/clock
datum} (Odijk's $\delta^s$, \cite{odijk2003fast}): the estimable
combination that, when removed by a user receiver, restores the integer
property of that user's single-difference ambiguities. This lumping is
exactly why a \emph{consistent} phase-bias product matters
(\dcref{sec:ssr}{the SSR-corrections section}): the producer's datum
and the consumer's integer recovery must be the same S-system.
\end{remark}

\subsection{Rank deficiency and the minimum-constraint S-basis}
\label{sec:net-sbasis}

The undifferenced network system \Cref{eq:net-station-rows} is rank
deficient. The same linear dependencies that the user side disposes of
by differencing are present here as exact \emph{gauge freedoms} of the
joint normal matrix: there is no observation that distinguishes a common
shift of all satellite clocks from an opposite shift of all receiver
clocks, nor the split of an absolute phase bias between the satellite
column and the per-station ambiguity datum, nor an inter-system bias
(ISB) common to a constellation. These are precisely the S-system
ambiguities of \cite{teunissen1985generalised}: the data fix only
\emph{estimable combinations}, and a \emph{datum} must be chosen to pin
the remaining null space.

\paragraph{Exact gauge, soft elimination.}
The chosen resolution is a \emph{minimum-constraint} datum applied as
exact hard pins to the reduced shared system, never a hard structural
elimination of columns. Concretely, each gauge is realised as a single
\emph{Dirichlet pin} of one shared column $c$ to a datum value (almost
always $0$): the column's known contribution is moved to the
right-hand side, its row and column in the reduced normal matrix are
zeroed, and a diagonal entry is set so that $x_c = \text{value}$ holds
exactly (\texttt{NetworkEstimator::solve},
\texttt{src/core/filter/NetworkEstimator.cpp}). The pin is applied to a
\emph{copy} of the reduced normals (the kernel's
\texttt{solveSharedConstrained} seam), leaving the estimator's own
reduced system pristine so per-station back-substitution
(\Cref{sec:net-schur}) stays valid. The pinned diagonal is set
commensurate with the reduced normals' own scale rather than to unity,
so that a genuinely \emph{under}-specified datum still surfaces as a
near-zero \emph{relative} pivot in the positive-definiteness guard,
independent of the network's overall weighting magnitude.

In S-transform language, if $\bm{x}$ collects the rank-deficient
undifferenced parameters and $\bm{S}$ is the S-basis selection, the
emitted satellite corrections are the estimable combination
$\bm{S}\,\bm{x}$; choosing the datum \emph{is} choosing $\bm{S}$, and
the hard pins above are the minimal set of rows of $\bm{S}$ that remove
the deficiency. The two halves of the datum live in two places:

\begin{itemize}[leftmargin=*, itemsep=1pt]
\item \textbf{Shared-block gauges} --- the overall clock datum (one per
      constellation), the per-constellation phase-bias$\leftrightarrow$ambiguity
      datum, and the ISB datum --- are the elements that touch the shared
      satellite columns. These go through the explicit
      \texttt{NetworkEstimator::Datum} (a list of column pins) and are
      applied in \texttt{solve}.
\item \textbf{Local datum elements} --- a network \emph{pivot
      receiver}'s own ambiguity/clock pin, which the emitted satellite
      phase bias absorbs (Odijk Eq.~4/9) --- live on that one station's
      private block. They are realised by the \emph{caller} as tight
      prior rows folded into that station's $\bm{A}_{r,\rx}$, $\bm{W}$,
      and $\bm{y}$, not through the shared \texttt{Datum}
      (\texttt{NetworkEstimator.hpp}). Station priors are also what make
      each local block $\Nrr$ full rank for the Schur step
      (\Cref{sec:net-schur}).
\end{itemize}

\paragraph{Selectable datum realisations.}
The pins just described are the \emph{legacy default}; each half of the
datum now admits alternatives, selected by config and recorded in the
\texttt{SbasisDeclaration}, all producing the same estimable
$\bm{S}\,\bm{x}$ but distributing the gauge differently. The clock/ISB
datum offers three modes (\texttt{SbasisDeclaration::IsbDatum}): the
hard \texttt{PerConstellationSatClockPin} above (default), a
\texttt{SumToZeroPerConstellation} constraint
(\texttt{network.clock\_datum\_sum\_to\_zero}) that spreads the pin
across all satellites of a constellation, and a
\texttt{BroadcastAnchoredPerSatellite} soft anchor
(\texttt{network.clock\_datum\_broadcast\_anchored}) tying each
satellite clock to its broadcast value at a declared $1\sigma$. The
ambiguity$\leftrightarrow$phase-bias datum likewise offers three
(\texttt{AmbDatum}): the local \texttt{PivotReceiverZeroAmbiguity} pin
above (default), a \texttt{SumToZeroPerReceiver} ensemble
(\texttt{network.phase\_datum\_ensemble}), and a
\texttt{SingleReceiverChain} realisation
(\texttt{network.phase\_datum\_eq34}). When satellite code biases are
estimated (the W6 slice), their own gauge --- a per-constellation
sum-to-zero (\texttt{CodeBiasDatum::SumToZeroPerConstellation}) --- is
folded as explicit constraint rows through
\texttt{NetworkEstimator::addSharedConstraintRows}, not as column pins.
Whichever realisations are active, the rank gate below is declared
against the matching S-basis dimension.

\begin{remark}
A rigorous null-space rank gate --- eigen-decompose the pre-constraint
reduced normals and assert that the count of small eigenvalues equals
the \emph{declared} S-basis dimension --- is implemented as
\texttt{NetworkSbasis::assertNullSpaceRank} and is invoked by the driver
on \texttt{reducedSharedNormals()} \emph{before} the constrained solve,
against the \texttt{SbasisDeclaration}'s \texttt{declared\_dimension}
(\texttt{Application::runNetwork}). It complements the
positive-definiteness guard: the guard catches the exact-deficiency
(under-specified datum) case on the pinned pivot, while the rank gate
certifies the reduced deficiency against the declared gauge dimension
and withholds any epoch whose null-space count disagrees --- re-gating
once after a drop-and-heal pass on the offending near-null direction.
\end{remark}

\subsection{Per-station Schur reduction (arrowhead normals)}
\label{sec:net-schur}

Stacking \Cref{eq:net-station-rows} over all stations and forming
weighted normal equations $\bm{A}^\Tr \bm{W} \bm{A}\,\bm{x} =
\bm{A}^\Tr \bm{W} \bm{y}$ yields a \emph{bordered block-arrow}
(``arrowhead'') system, because the station--station off-diagonal
blocks are exactly zero (stations couple only through $\xshared$):
\begin{equation}
\begin{bmatrix}
\Nss      & \bm{N}_{s1} & \bm{N}_{s2} & \cdots & \bm{N}_{sR} \\
\bm{N}_{1s} & \bm{N}_{11} & \bm{0}      & \cdots & \bm{0} \\
\bm{N}_{2s} & \bm{0}      & \bm{N}_{22} & \cdots & \bm{0} \\
\vdots      &             &             & \ddots &        \\
\bm{N}_{Rs} & \bm{0}      & \bm{0}      & \cdots & \bm{N}_{RR}
\end{bmatrix}
\begin{bmatrix}
\xshared \\ \xlocal{1} \\ \xlocal{2} \\ \vdots \\ \xlocal{R}
\end{bmatrix}
=
\begin{bmatrix}
\bm{u}_s \\ \bm{u}_1 \\ \bm{u}_2 \\ \vdots \\ \bm{u}_R
\end{bmatrix},
\label{eq:net-arrowhead}
\end{equation}
with the per-station blocks built from that station's pre-formed design
rows,
\begin{equation}
\Nss^{\,\rx} = \bm{A}_{s,\rx}^\Tr \bm{W}_\rx \bm{A}_{s,\rx},
\quad
\Nsr^{\,\rx} = \bm{A}_{s,\rx}^\Tr \bm{W}_\rx \bm{A}_{r,\rx},
\quad
\Nrr^{\,\rx} = \bm{A}_{r,\rx}^\Tr \bm{W}_\rx \bm{A}_{r,\rx},
\label{eq:net-blocks}
\end{equation}
and $\bm{u}_s^{\,\rx} = \bm{A}_{s,\rx}^\Tr \bm{W}_\rx \bm{y}_\rx$,
$\bm{u}_r^{\,\rx} = \bm{A}_{r,\rx}^\Tr \bm{W}_\rx \bm{y}_\rx$. A full
(non-diagonal) $\bm{W}_\rx$ is supported, so correlated undifferenced
measurement covariance (e.g.\ shared pre-corrections) propagates
exactly (\texttt{addStationObservations},
\texttt{NetworkNormalEquations.hpp}).

\paragraph{Lossless marginalisation.}
The arrowhead structure is solved by Schur-eliminating each station's
private block onto the shared border. Each $\Nrr^{\,\rx}$ is factorised
once ($\mathrm{LDL}^\Tr$) when the station is added; the elimination
accumulates
\begin{empheq}[box=\fbox]{align}
\redNss &= \sum_{\rx \in \rxset}
   \Big( \Nss^{\,\rx} - \Nsr^{\,\rx}\,(\Nrr^{\,\rx})^{-1}\,\Nsr^{\,\rx\Tr} \Big)
   \;+\; \bm{N}_{ss}^{\text{prior}},
\label{eq:net-schur-N}\\
\redus &= \sum_{\rx \in \rxset}
   \Big( \bm{u}_s^{\,\rx} - \Nsr^{\,\rx}\,(\Nrr^{\,\rx})^{-1}\,\bm{u}_r^{\,\rx} \Big)
   \;+\; \bm{u}_s^{\text{prior}},
\label{eq:net-schur-u}
\end{empheq}
the reduced shared system $\redNss\,\xshared = \redus$
(\texttt{NetworkNormalEquations::reduce}). The reduction is a
\emph{lossless} marginalisation, not an approximation: solving
\Cref{eq:net-schur-N,eq:net-schur-u} for $\xshared$ and
back-substituting per station,
\begin{equation}
\xlocal{\rx} = (\Nrr^{\,\rx})^{-1}\big(\bm{u}_r^{\,\rx} - \Nsr^{\,\rx\Tr}\,\xshared\big),
\label{eq:net-backsub}
\end{equation}
reproduces the dense stacked solve to machine precision. Fill-in is
confined to the satellite border and each station's elimination is
independent, so the cost scales linearly in the station count and the
per-station eliminations parallelise (\texttt{reconstructStation},
\texttt{NetworkNormalEquations.cpp}). The shared-block priors
$\bm{N}_{ss}^{\text{prior}}$ in \Cref{eq:net-schur-N} --- soft physical
priors and any additive exact-gauge contributions --- are added
pre-reduction via \texttt{addSharedInfo} and carry through the Schur
complement unchanged.

\begin{remark}
This arrowhead-with-exact-reduction topology mirrors the NGS M-PAGES
\texttt{NgsLsq::NormalEquations::reduce}\,/\,\texttt{reconstructLocals}
batch network solver, specialised so that the M-PAGES ``global'' block
is the shared satellite border and the per-epoch swept block becomes the
$R$ independent station blocks. The design panel rejected a two-stage
covariance-form alternative precisely because the normal-equation Schur
form is exact and cheaper here.
\end{remark}

\subsection{Constellation scope}\label{sec:net-constellations}

The producer estimates GPS and Galileo, with BeiDou supported on the
same footing; \textbf{GLONASS phase biases are excluded}. The reason is
structural, not a tuning choice: GLONASS is FDMA, so each satellite
transmits on a slightly different frequency and the receiver
inter-channel (hardware) bias does \emph{not} cancel cleanly across
satellites the way it does for the CDMA constellations. The
satellite/receiver phase-bias bookkeeping that underpins the
explicit-bias S-system
(\Cref{sec:net-phasebias,sec:net-sbasis}) assumes a per-signal bias
common across a constellation's satellites; the per-frequency GLONASS
inter-channel bias breaks that assumption and would corrupt the phase
datum. GLONASS is therefore excluded from the network \emph{phase}-bias
product by design
(\texttt{docs/ssystem\_conformance\_2026-06-15.md}); the user side
excludes it from ambiguity resolution for the same reason. A
\emph{code-only} GLONASS mode is nonetheless supported behind
\texttt{network.estimate\_glonass} (\texttt{enableGlonassNetwork},
default off): it admits GLONASS to the solve carrying a satellite clock
and a per-slot receiver inter-frequency bias but \emph{no} FDMA phase
bias, so the phase-datum bookkeeping above is untouched. (This is
distinct from the float-only GLONASS support the user engine carries for
positioning; with the code-only mode off, the network \emph{producer}
generates no GLONASS bias products at all.)

\subsection{The closed loop: producer $\to$ held-out user}
\label{sec:net-closedloop}

The emitted product is one \texttt{NetworkSolutionEpoch} per content
epoch: per-satellite clock $C_0$ plus per-signal code and phase biases,
keyed by RINEX satellite id and serialised in IGS-SSR convention by
\texttt{NetworkProductEmitter}
(\texttt{NetworkEstimator::toSolutionEpoch}). That product is exactly
the SSR stream the single-receiver engine already knows how to consume:
the satellite clock refines $-\sLight\,\mathrm{d}t^\sat$ and the
observable-specific biases refine $-\satbiasCB^\sat_\sigl$ and
$-b^\sat_\sigl$ as \emph{known} corrections, restoring the integer
ambiguity property exactly as described in
\dcref{sec:ssr}{the SSR-corrections section}.

The validation strategy follows directly from this symmetry. The
companion \texttt{ProcessingMode::PPP\_RTK\_User} mode runs the ordinary
uncombined PPP-RTK pipeline but consumes the network's own
bias/atmosphere product. Holding one station \emph{out} of the network
solve and then positioning it as a user against the produced
corrections turns the consumer into a closed-loop product test: the
held-out station's recovered position and fixed ambiguities certify that
the producer's S-system datum and the consumer's integer recovery are
mutually consistent. The \texttt{PPP\_RTK\_User} consumer flows through
the standard PPP path (it is not a separate estimator); only the
producer takes the joint multi-station pipeline of this section.

\ifSubfilesClassLoaded{\printbibliography}{}
\end{document}
